\section{超声波振动的减阻特性试验}

\subsection{试验方法}

施肥机在排肥作业过程中，有机肥与肥箱底部钢板之间存在大量相对运动，存在较大的接触压力，由于有机肥颗粒与钢板的摩擦系数的影响，施肥机在排肥作业过程中排肥装置受到较大阻力，因此通过探究超声振动对有机肥摩擦系数的影响规律从而揭示超声振动可以改变有机肥的流动特性。

本次试验探究施加超声振动对有机肥摩擦系数的影响，需要得到准确的摩擦系数，试验物料与材料接触表面间的静摩擦系数一般可通过斜面法和拖重法两种方法测量，粉末颗粒都选用斜面法测量摩擦系数。因此本试验运用自制试验装置（图\ref{pic:斜面仪法}）测量超声振动下的有机肥摩擦系数，通过自吸式数显角度仪对有机肥颗粒和钢板材料的静摩擦角进行测定。

试验时，将数显角度仪吸附斜面钢板固定位置处，应用数显角度仪调节斜面装置水平度，将70mg有机肥样品放置于钢板一侧的有机肥放置位置处，换能器通过环氧树脂固定安装于斜面钢板另一侧表面中心。逆时针缓慢转动斜面钢板至有机肥颗粒在钢板开始滑动时停止转动，记录该时刻数显角度仪器的角度$\theta_1$，计算静摩擦系数值$\mu_1$。

\begin{equation}
	\mu_1 = \tan {\theta_1} \label{eq:静摩擦系数}
\end{equation}

\begin{figure}[H]
	\centering
	\begin{tikzpicture}
	\begin{scope}[rotate around={30:(0,0)}]
		\draw[fill=white] (0,0) rectangle ($(5,0) + (0,0.1)$);
		\draw[fill=black] ($(2.5,0)+(0.5,0.1)$) arc (0:180:0.5);
		\draw[fill=gray] (5,0) -- ($(2.5,0)+(0.25,0)$) -- ($(2.5,-0.5)+(0.15,0)$) -- ($(2.5,-0.5)-(0.15,0)$) -- ($(2.5,0)-(0.25,0)$) -- cycle;
	\end{scope}
	\draw[thick] (0,0) to (5,0);
	\foreach \x in {0,0.2,...,5}{
		\draw (\x,0) -- ({\x+0.5},-0.5);
	}
	\draw[thick] ($(0,0)+(1,0)$) arc (0:30:1);
	\coordinate (A) at (15:1.2);
	\node[anchor=center] at (A) {$\theta$};
	\coordinate (B) at (1.8,1.8);
	\draw (B) -- +(-0.5,0.5) node[above left]{1};
%	\draw[color=gray,step=0.2] (0,0) grid (5,5);
%	\draw[color=red,step=1] (0,0) grid (5,5);
	\draw (2.6,1) -- +(0.5,0) node[right]{2};
	\draw (3.8,2.4) -- +(-0.5,0.2) node[above left]{3};
	\draw[->] (5:5) arc (5:25:5);
\end{tikzpicture}


	{\small 1-有机肥堆；2-超声波换能器；3-钢板；}
	\caption{斜面仪法}
	\label{pic:斜面仪法}
\end{figure}

为了便于比较，定义超声波振动减摩率如式\eqref{eq:减摩率}所示。

\begin{equation}
	\alpha = \frac{\tan{\theta_0}-\tan{\theta_1}}{\tan{\theta_0}} \label{eq:减摩率}
\end{equation}

式中，$\alpha$为超声波振动减摩率；$\theta_0$为有机肥摩擦角；$\theta_1$为超声波振动感激励下的有机肥摩擦角。

\subsection{用到的传感器}

该试验主要使用了倾斜角传感器，其主要工作部分是加速度计（Accelerometer）。加速度计是一种测量加速度向量的传感器。它一般是MEMS（微机电）结构：质量块+弹簧+电容板。

\begin{figure}[H]
	\centering
	\includegraphics[width=0.5\textwidth]{}
\end{figure}

\subsection{试验方案}

\subsubsection{超声振动对有机肥摩擦系数影响单因素试验方案}

对超声振动对有机肥摩擦系数影响进行预试验,根据预试验结果筛选出超声振动  对有机肥摩擦系数影响显著的四个因素,分别为超声振动频率、超声振动功率、Q235  钢板厚度、有机肥含水率,以评价超声振动对有机肥摩擦系数影响,单因素试验方案如  下:

\begin{enumerate}
	\item 超声功率对有机肥摩擦角的影响

	考虑到换能器最大输入功率为60W，因此本研究控制超声振动功率为10-60W，期间每10W为一个梯度，共六个功率水平，控制超声振动频率为28kHz，Q235钢板厚度3mm，有机肥含水率35\%。研究不同超声功率对有机肥摩擦系数的影响，每组试验重复5次，取其各自测量的平均值作为最终测量的结果。为确保有机肥含水率的准确性，每次试验更换样品。

	\item 超声频率对有机肥摩擦角的影响
	
	目前超声振动换能器主要有25kHz，28kHz和40kHz，因此本研究控制超声振动频率为25kHz、28kHz和40kHz，共三个频率水平，控制超声振动功率为35W，Q235钢板厚度3mm，有机肥含水率35\%。研究不同超声频率对有机肥摩擦系数的影响，每组试验重复5次。

	\item Q235钢板厚度对有机肥摩擦角的影响
	
	由于肥箱底部钢板厚度一般为2mm左右，因此本研究控制Q235钢板厚度为15mm，每1mm为一个梯度，共五个水平。控制超声振动频率为28kHz，超声振动功率为35W，有机肥含水率35\%。研究不同Q235钢板厚度对有机肥摩擦系数的影响，每组试验重复5次。

	\item 有机肥含水率对有机肥摩擦角的影响
	
	由于有机肥含水率范围较广，容易吸收空气中水分，因此本研究控制有机肥含水率为10\%-60\%，每10\%为一个梯度，共六个水平。控制超声振动频率为28kHz，超声振动功率为35W，钢板厚度3mm。研究不同有机肥含水率对有机肥摩擦系数的影响，每组试验重复5次。


\end{enumerate}

\subsubsection{超声振动对有机肥摩擦系数影响正交试验方案}

为了进一步探究不同因素交互作用下对有机肥摩擦系数的影响，在单因素试验基础上，选取超声振动频率、超声振动功率和有机肥含水率三个因素以分析三因素交互作用下对有机肥摩擦系数的影响。正交试验的3个因素名称分别以符号正交试验的3个因素名称分别以符号A、B、C表示，各因素分别被赋予编码值dddddddddddddd1、0、-1以表示其水平的差异，具体试验水平表如表所示。正交试验基于Design Expert 12软件进行，设置3个中心点用于误差分析，共进行15组试验，具体试验方案与结果见表\ref{tb:正交因素试验和水平}。

\begin{table}[H]
	\centering
	\caption{正交因素试验和水平}
	\label{tb:正交因素试验和水平}
	\begin{tabularx}{\textwidth}{
		>{\centering\arraybackslash\hsize=1.0\hsize\linewidth=\hsize}X
		>{\centering\arraybackslash\hsize=0.5\hsize\linewidth=\hsize}X
		>{\centering\arraybackslash\hsize=0.5\hsize\linewidth=\hsize}X
		>{\centering\arraybackslash\hsize=0.5\hsize\linewidth=\hsize}X
		>{\centering\arraybackslash\hsize=0.5\hsize\linewidth=\hsize}X
		}
		\toprule
		\multirow{2}{*}{试验因素} & \multirow{2}{*}{编码值} & \multicolumn{3}{c}{试验水平} \\
		\cmidrule{3-5}
		 & & -1 & 0 & 1 \\
		\midrule
		超声波振动功率（W） & A & 10 & 35 & 60 \\
		超声波振动频率（kHz） & B & 25 & 28 & 40 \\
		有机肥含水率（\%） & C & 10 & 35 & 60 \\
		\bottomrule
	\end{tabularx}
\end{table}